\section{Propiedades y Desafíos del Aprendizaje sobre Grafos}

\subsection{Aprendizaje Inductivo frente a Transductivo}

La distinción entre aprendizaje transductivo e inductivo resulta central para comprender qué puede hacer un modelo de grafos una vez desplegado. En el paradigma transductivo, el modelo observa durante el entrenamiento la totalidad del grafo, incluidos los nodos cuyas etiquetas se desean predecir, de modo que aprende representaciones ligadas a posiciones concretas de esa estructura fija. La formulación original de la \textit{Graph Convolutional Network} pertenece a esta categoría, porque sus pesos se ajustan sobre una matriz de adyacencia conocida de antemano y no ofrece un mecanismo natural para incorporar vértices ausentes en el momento del ajuste. Esta característica limita su aplicabilidad cuando el objetivo consiste en clasificar entidades que todavía no existían al entrenar, ya que cualquier incorporación obligaría en principio a repetir el proceso completo sobre el grafo ampliado.

El enfoque inductivo supera esa restricción al aprender funciones de agregación que son transferibles a cualquier vecindario, en lugar de memorizar un embedding por nodo. Arquitecturas como \textit{GraphSAGE} y la \textit{Graph Attention Network} encarnan esta idea, puesto que entrenan operadores que combinan los atributos de un nodo con los de sus vecinos siguiendo reglas compartidas por todo el grafo. Gracias a ello, cuando aparece un vértice nuevo basta con recorrer su vecindario local y aplicar las mismas transformaciones aprendidas para obtener su representación, sin necesidad de reentrenar. Esta transferibilidad convierte al modelo en una función generalizable sobre la estructura, y no en una tabla de valores atada a un conjunto cerrado de identidades.

La capacidad inductiva adquiere una importancia decisiva en el dominio de las redes financieras, donde el grafo es intrínsecamente dinámico y crece de forma continua. Cada día se abren cuentas y se registran transacciones que introducen vértices y aristas antes inexistentes, de manera que un sistema de detección obligado a reentrenar ante cada llegada sería inviable en la práctica. Un modelo inductivo, en cambio, evalúa una operación recién observada apoyándose en su contexto local inmediato, lo que se vincula de forma directa con las técnicas de muestreo de vecindarios. Dicho muestreo selecciona un subconjunto acotado de vecinos en cada capa para mantener el cómputo tratable sobre grafos masivos, y a la vez sostiene la premisa de que la representación de un nodo depende de su entorno próximo y no de la totalidad del grafo, lo que refuerza la coherencia entre eficiencia y generalización.

\subsection{El Sobre-suavizado y el Límite de la Profundidad}

El fenómeno conocido como sobre-suavizado, u \textit{over-smoothing}, describe una degradación progresiva que sufren las representaciones de nodos a medida que se apilan capas de paso de mensajes. En cada capa, un nodo actualiza su representación agregando información de sus vecinos, de modo que iterar esta operación equivale a promediar señales sobre vecindarios cada vez más amplios. Cuando el número de capas crece demasiado, las representaciones de nodos que originalmente eran distintos tienden a converger hacia un valor común, y el grafo entero se aproxima a un estado en el que casi toda la información local se ha diluido. El resultado es una pérdida de poder discriminativo, porque el clasificador situado al final recibe vectores que apenas difieren entre clases y por tanto separa mal las categorías de interés.

Esta dinámica explica por qué la profundidad útil de una \textit{Graph Neural Network} suele quedar acotada a unas pocas capas, en marcado contraste con las redes profundas de otros dominios como la visión por computador, donde decenas o cientos de capas mejoran el desempeño. La diferencia radica en que apilar capas en un grafo no solo aumenta la abstracción, sino que además expande el radio del vecindario del que cada nodo extrae información, y ese doble efecto genera una tensión difícil de conciliar. Por un lado interesa ampliar el campo receptivo para capturar dependencias entre nodos distantes que puedan ser relevantes; por otro lado, ampliarlo en exceso arrasa con las diferencias locales que permiten distinguir un nodo de otro. El diseño de la arquitectura debe entonces buscar un punto de equilibrio, en lugar de suponer que más profundidad siempre aporta mayor capacidad.

Las consecuencias prácticas de este límite se manifiestan al momento de decidir cuántas capas emplear, decisión que deja de ser un detalle menor para convertirse en una elección de fondo. Un número reducido de capas preserva la nitidez de las representaciones, pero restringe el alcance del contexto que el modelo puede integrar, con lo que patrones que se despliegan a lo largo de trayectorias largas pueden quedar fuera de vista. Un número elevado amplía ese alcance a costa de acercarse al régimen de sobre-suavizado, donde la ventaja teórica del contexto extendido se desvanece por la homogeneización de las señales. Por ello, la selección del número de capas se plantea como un compromiso entre alcance y distinción, ajustado a la escala de las estructuras que se pretende reconocer en cada problema concreto.

\subsection{Dinámica de Entrenamiento y Regularización en GNN}

El ajuste de una \textit{Graph Neural Network} se realiza mediante descenso de gradiente, procedimiento que modifica de forma iterativa los pesos para reducir una función de pérdida que cuantifica el error de predicción. El cálculo de los gradientes se apoya en la retropropagación, o \textit{backpropagation}, que propaga la señal de error desde la salida hacia las capas anteriores aplicando la regla de la cadena a lo largo del grafo de cómputo. Sobre esa base, el optimizador \textit{Adam} suele encargarse de gobernar las actualizaciones, puesto que adapta el tamaño del paso para cada parámetro combinando estimaciones del primer y del segundo momento del gradiente, lo que estabiliza el avance en superficies de error irregulares. La tasa de aprendizaje regula la magnitud global de cada paso, y el decaimiento de pesos añade una penalización sobre la norma de los parámetros que desalienta soluciones de magnitud excesiva y favorece una mejor generalización.

Para contener el \textit{overfitting}, es decir, la tendencia del modelo a memorizar el conjunto de entrenamiento en lugar de capturar regularidades transferibles, se recurre habitualmente al \textit{dropout}. Esta técnica desactiva de manera aleatoria una fracción de las unidades durante el entrenamiento, de modo que la red no puede depender en exceso de neuronas individuales y se ve forzada a distribuir la información entre múltiples rutas. El efecto se asemeja a entrenar de forma implícita un conjunto de subredes que luego se combinan, lo que reduce la varianza del modelo y mejora su comportamiento sobre datos no vistos. En conjunto, la tasa de aprendizaje, el decaimiento de pesos y el \textit{dropout} conforman un repertorio de mecanismos de regularización cuyo ajuste conjunto condiciona de manera notable la calidad final del modelo.

El entrenamiento sobre grafos plantea, sin embargo, desafíos que lo apartan del caso habitual de datos independientes. Las muestras no son autónomas entre sí, porque la predicción de un nodo depende de sus vecinos, de manera que la noción de \textit{batch} debe redefinirse en términos de subgrafos que preservan la conectividad local necesaria para el paso de mensajes. Esta dependencia complica la partición de los datos y exige cuidado para evitar que información del conjunto de evaluación se filtre a través de las aristas hacia el entrenamiento. A ello se suma la sensibilidad a la inicialización de los pesos y a la aleatoriedad del muestreo, que puede llevar a que ejecuciones distintas converjan a soluciones apreciablemente diferentes, lo que motiva repetir el entrenamiento con varias semillas para caracterizar la variabilidad y sostener conclusiones más robustas.

\subsection{Capacidad Expresiva de las Redes Neuronales de Grafos}

El poder expresivo de una \textit{Graph Neural Network} se refiere a qué estructuras del grafo es capaz de distinguir a través de las representaciones que produce. Una arquitectura más expresiva puede asignar embeddings diferentes a subestructuras genuinamente distintas, mientras que una menos expresiva las confunde al proyectarlas sobre representaciones idénticas. Esta noción importa porque el valor de un modelo no reside solo en su exactitud sobre un conjunto particular, sino en su capacidad teórica de separar configuraciones que difieren en su topología. Comprender ese techo resulta esencial para anticipar qué clases de patrones quedan, por principio, fuera del alcance de una familia de modelos.

Buena parte de las arquitecturas de paso de mensajes encuentran su límite superior de expresividad en el test de isomorfismo de grafos de Weisfeiler-Lehman en su variante de una dimensión. Dicho test refina de forma iterativa una etiqueta por nodo agregando las etiquetas de sus vecinos, exactamente el mismo esquema de vecindad que emplea el paso de mensajes, razón por la cual ambos comparten la misma frontera de discriminación. La consecuencia es que si el test no logra diferenciar dos grafos, tampoco lo hará una red que opere bajo ese esquema, con independencia de cuántos parámetros o capas incorpore. Esta correspondencia ofrece un marco preciso para razonar sobre las capacidades y las limitaciones estructurales de las arquitecturas más difundidas, y sitúa un límite claro que el diseño por sí solo no puede rebasar.

De este límite se desprende que dos grafos distintos pueden generar el mismo embedding cuando sus vecindarios resultan indistinguibles para el mecanismo de agregación, aun cuando difieran en propiedades globales relevantes. Para la detección de lavado de dinero, esta limitación adquiere un peso considerable, ya que ciertos esquemas ilícitos se apoyan precisamente en configuraciones estructurales sutiles concebidas para diluirse entre la actividad legítima. Patrones como ciclos que retornan fondos a su origen o topologías de estratificación diseñadas para fragmentar rastros pueden quedar sin diferenciar si su firma estructural coincide, a ojos del modelo, con la de operaciones normales. Reconocer esta cota orienta el trabajo hacia el enriquecimiento de los atributos de nodos y aristas o hacia arquitecturas de mayor expresividad, y a la vez modera las expectativas sobre lo que un modelo dado puede captar apoyándose únicamente en la estructura del grafo.
